\section{Experimental Design}\label{sec:design}

\subsection{Benchmark, evaluator, and evidence tiers}
RoboTwin 2.0 is the primary benchmark.  We evaluate hammering, ranking by
color, ranking by size, handover, stacking two blocks, and stacking three
blocks.  Each completed pilot cell contains four simulator seed groups and 50
accepted expert-feasible episodes per group (200 episodes per task and 1,200
per method).  Policies predict 50 actions and replan after executing 36.
Accepted scene identifiers must be intersected across conditions for paired
tests.  LIBERO is reported only as an auxiliary saturation and negative-transfer
diagnostic.

We separate three evidence tiers.  The \emph{external calibration} is the
published $\pi_{0.5}$ checkpoint evaluated through our bridge.  The
\emph{legacy pilot} starts every arm from \texttt{pi05\_base} but uses
joint-delta actions, quantile normalization, batch 64, and 20k updates.  The
\emph{confirmatory matrix} will use absolute actions, mean/std normalization,
batch 16, 50k updates, and the published optimizer schedule for every arm.
Only the confirmatory matrix can support the final Q1 claim.

\subsection{Q1: matched control performance}
The claim-bearing comparison is A0 (no hint), A2-Abs (offline absolute
milestone), and A3/MINT-VLA (current-encoder residual predictor), trained from
the same initialization and data under the confirmatory recipe.  A0 has no
hint field, milestone loss, or milestone-specific parameters.  We will report
each training seed, the seed mean and interval, and all task-level results.

\exptodo{T0--T2}{Finish corrected A0; after its 70\% operational gate, train
matched A2-Abs and six-task-coverage A3 for three independent seeds.}

\subsection{Q2: effect of the predicted token content}
For each trained A2-Abs and A3 checkpoint, we hold weights fixed and replace the
correct prediction with (i) the current-state feature, (ii) zeros, (iii) a
within-task shuffled future, and (iv) a cross-task future.  A completed
feature-dimension permutation is reported separately because it tests feature
structure, not instance relevance.  Conditions use the
same accepted scene identifiers.  Correct-versus-current tests temporal content;
correct-versus-shuffled tests instance relevance while preserving a plausible
feature distribution; zero tests removal of the token; cross-task tests semantic
misdirection.  We report paired success differences and exact McNemar tests
with Holm correction across planned contrasts.
\exptodo{E1}{Complete A3 zero, instance-shuffled, and cross-task conditions;
run the missing A2 instance shuffle; freeze scene intersections and report
discordant counts with Holm-adjusted exact McNemar tests.}

\subsection{Q3: factorized interface mechanism}
The A2-to-A3 pilot changes how target features are computed, predictor capacity, auxiliary
loss, and the injected value at once.  The confirmatory ablation
therefore holds predictor architecture, loss, injected value, and parameter
count fixed while varying one factor:
\begin{center}
\small
\begin{tabular}{p{0.31\linewidth}p{0.55\linewidth}}
\toprule
Factor & Controlled comparison \\
\midrule
target computation & saved target feature vs. frame re-encoded by the current encoder \\
predictor form & direct absolute predictor vs. residual-parameterized predictor, both injecting absolute \\
forward guidance & auxiliary-only/no-prefix vs. predicted prefix \\
backward shaping & prefix-only/no-aux and detached-current auxiliary branch \\
milestone selection & fixed horizon, random future, terminal frame, and oracle milestone \\
\bottomrule
\end{tabular}
\end{center}
\exptodo{T3}{Train the one-factor target-computation, predictor-form, forward/backward
route, capacity, and instruction-aware controls.}

\subsection{Q4: robustness, scope, and efficiency}
Training seed is the top-level uncertainty unit.  We will use a hierarchical
bootstrap over training seed, task, and paired scene, while retaining task-level
values to expose heterogeneous effects.  Transfer is tested by matched
fine-tuning from a mature initialization, not by comparing against an unrelated
public checkpoint.  The task-scope hypothesis is predeclared on block and bowl
stacking, with reactive/contact and fine-grained geometric tasks as controls.
We will also report added parameters, training FLOPs, peak memory, inference
latency, and throughput relative to A0.

Training-seed, mature-initialization, and predeclared task-group experiments
remain confirmatory gates.  We report static deployment cost now and reserve
measured memory and latency claims for a paired benchmark.
\exptodo{T4--T5,E2}{Train the predeclared task panel and mature-initialization
matrix; measure paired peak memory, inference latency, and server throughput.}

\subsection{Q5: instantiation on a second VLA architecture}
The milestone target and predictor are defined independently of the
$\pi_{0.5}$ prefix.  To test use with another architecture, we will retain the
milestone selection, native-space objective, data, and evaluation protocol;
size the predictor to the second VLA's visual feature width; and connect it
through that VLA's native conditioning input.  The identifying comparison is
MINT-VLA versus a matched no-milestone baseline within the second architecture,
not raw success differences between two VLA families.
\exptodo{T6}{Report matched no-milestone and MINT-VLA results for at least one
additional VLA; do not claim multi-architecture evidence from the $\pi_{0.5}$
experiment alone.}
